\scanpage{78}
essentielle triviale, qui fait \unsure{pendant|penser} \unsure{à|:} (6), est
\[
(2\unsure{8})\qquad {}^{uv}l={}^u({}^vl)
\qquad u,v\in\Gamma,\quad l:x\longrightarrow y
\quad\text{chemin dans }\unsure{\widehat{\Pi}_1(X)|\Pi_1(X)}
\]
\deleted{d'où,} si $l$ est chemin de $\Pi_1(X)$, d'où
\[
{}^vl=g_l(v)\cdot l
\]
\deleted{\unsure{grâce à} (6)} d'où on tire par (\unsure{6|8})
\[
{}^u({}^vl)={}^u\{g_l(v)\}\cdot \underset{\begin{gathered}
\rule{0.35pt}{3pt}\rule{2.2em}{0.35pt}\rule{0.35pt}{3pt}\\[-2pt]
g_l(u)\cdot l
\end{gathered}}{{}^ul}={}^ug_l(v)g_l(u)\cdot l
\]
d'où on tire
\[
(29)\qquad\left\{\begin{aligned}
{}^{\deleted{\unsure{u}}}g_l(uv)&={}^ug_l(v)\cdot g_l(u)\qquad\text{soit}\\
g_l^{-1}(uv)&=g_l^{-1}(u)\cdot{}^ug_l^{-1}(v)
\end{aligned}\right.
\]
Et de $=$
\[
(30)\qquad f_l(uv)=f_l(u)\cdot{}^uf_l(v)
\]
i.e. $f_{l\alpha}$ et $g_l^{-1}$ sont, pour $l$\unsure{.} fixé, $u$ variable, des $1$-cocycles \inserted{sur $\Gamma$} à valeurs dans $\widehat{\Pi}_1(X,x)$ resp. $\widehat{\Pi}_1(X,y)$, \unsure{\inserted{(groupes sur lesquels $\Gamma$ opère)}|(groupes sur lesquels $\Gamma$ opère)}, d'ailleurs déduits l'un de l'autre par l'isomorphisme
\diagram{\begin{tikzcd}[ampersand replacement=\&]
\widehat{\Pi}_1(X,x)
\arrow[r,phantom,"{\begin{gathered}\Pi_1(l)\\[-2pt]\unsure{\relbar\joinrel\relbar|\longrightarrow}\end{gathered}}"]
\&\widehat{\unsure{\Pi|\pi}}_1(X,y)
\end{tikzcd}}
induit par $l$ \deleted{\unsure{en $l$}} Ici, c'est la formule pour $f_l$ qui a une gueule plus sympathique que celle pour $g_l$ (pour (10)(11) c'était l'inverse --- \inserted{exprim.} ``multiplicativité en $l$''). Les $1$-cocycles obtenus (\deleted{\unsure{et}} $f$ \inserted{un $1$-cocycle} droit, $g$ un $1$-cocycle à gauche) correspondent à des $\Gamma$ torseurs à droite resp. à gauche sous $\widehat{\Pi}_1(X,x)$ resp! $\widehat{\unsure{\pi|\Pi}}_1(X,y)$, qui bien sûr n'est autre que $\operatorname{Hom}_{\widehat{\Pi}_1(X)}(x,y\unsure{)}$

\margin{On pourrait expliciter \unsure{un de vue qui voit le} couple $(f_l,g_l)$ comme un ``bicocycle'' sur $\Gamma$, relatif à un $\Gamma$-bitorseur $\operatorname{Hom}_{\widehat{\Pi}_1(X)}(x,y)$ sous $(\widehat{\Pi}_1(X,y),\widehat{\Pi}_1(X,x))$}
\margin{77}
